% sec_introduction.tex — Section 1: Introduction (compressed)
% ICGNC 2026, Springer svproc template

\section{Introduction}
\label{sec:intro}

Multi-UAV cooperative search and attack missions require a swarm of UAVs to
autonomously explore unknown environments, discover hidden targets, and
coordinate strikes under stringent communication constraints and threat
zones~\cite{shakhatreh2019unmanned}.  The core challenge is achieving effective
inter-agent coordination when each UAV possesses only a partial, local view and
can communicate only within a limited radius.

Multi-Agent Reinforcement Learning (MARL) offers a data-driven approach to
learning cooperative policies~\cite{lowe2017multi,yu2022surprising}.
Graph-based communication mechanisms such as
CommNet~\cite{sukhbaatar2016learning} and DGN~\cite{jiang2020graph} enable
learned message passing, yet standard Graph Attention Networks
(GATs)~\cite{velickovic2018gat} treat all links equally, ignoring that wireless
channel quality degrades with distance, co-modal agents share higher mutual
information, and stale messages carry increasing uncertainty.  Furthermore, MARL
policies are inherently opaque.  Large Language Models
(LLMs)~\cite{brown2020language} have shown strong reasoning
capabilities~\cite{yao2023react,shinn2023reflexion}, but directly applying them
to multi-UAV coordination is hindered by the lack of domain-specific
cooperative strategies and the difficulty of encoding topology-dependent states
into prompts.

We propose TAGA-LLM, a two-stage framework that combines MARL's
cooperative learning strength with LLM interpretability.  The key insight is
that MARL discovers effective cooperative strategies but locks them inside
opaque neural networks, while LLMs excel at structured reasoning and natural
language explanation but lack domain-specific cooperative knowledge;
distillation bridges this complementary gap.  In \emph{Stage~I}, we introduce the
\emph{Topology-Adaptive Graph Attention} (TAGA) mechanism with three learnable
factors---distance decay, task awareness, and information freshness---integrated
into a hierarchical MAPPO framework with curriculum learning.  In
\emph{Stage~II}, we distill the TAGA-MAPPO policy into a Qwen-4B LLM via LoRA
fine-tuning with topology-aware prompt encoding, CoT supervision, and
counterfactual data augmentation.

The primary contributions of this paper are summarized as follows:
\begin{inparaenum}[(a)]
\item a TAGA mechanism with three physics-informed adaptive factors for
  communication aware multi-agent coordination;
\item a two-stage MARL to LLM distillation pipeline that yields
  interpretable and instruction-controllable mission planning; and
\item a hierarchical policy architecture trained with MAPPO and a
  three-stage curriculum.
\end{inparaenum}

The rest of this paper is organized as follows.
In Sect.~\ref{sec:related}, we review related work.
In Sect.~\ref{sec:problem}, we state the Dec-POMDP formulation of the mission.
In Sect.~\ref{sec:method}, we present the TAGA-LLM framework, including TAGA,
MAPPO training, and distillation.
In Sect.~\ref{sec:experiments}, we report simulation and evaluation plans.
Finally, Sect.~\ref{sec:conclusion} concludes the paper.
